\section{Discussion}
\subsection{PA1 Performance and Dimensionality}
The imbalance explains why accuracy must be considered alongside minority-class metrics. The full-feature SVM combines accuracy 0.8653 and precision 0.8363 with recall 0.2241: its positive predictions are relatively precise, but it identifies only a small proportion of employees leaving. The full-feature Neural Network offers the strongest observed balance according to F1 and balanced accuracy, although recall 0.4219 still indicates substantial missed attrition cases.

Selection does not universally improve performance. K-NN benefits from univariate and wrapper selection, with the wrapper yielding its largest F1 improvement; multivariate selection instead lowers its F1. Removing dimensions can change which observations are considered neighbours, so retaining a useful subset can help this classifier. This interpretation is consistent with the observed improvements but does not establish the mechanism. Decision Tree improves moderately with the filters. All reduced Neural Network configurations lose performance relative to its complete representation, suggesting that the omitted information is useful in that experimental setting.

The filters assess feature--target association, with the multivariate filter also penalising redundancy. For example, the univariate subset retains both complementary overtime indicators, whereas the multivariate subset retains only \texttt{OverTime\_Yes}. The wrapper evaluates candidate additions through a Decision Tree's F1, so the resulting representation is linked to that estimator rather than independently optimised for every final classifier. These different criteria explain why the methods need not choose identical subsets or favour the same classifier. The 70.59\% reduction offers a more compact representation, but 15 transformed features do not necessarily represent 15 distinct employee characteristics. No selection strategy is best on every metric.

\subsection{Integrated Performance and Probability Quality}
Full-feature Logistic Regression in PA2 exceeds PA1's full-feature Neural Network in observed F1, while their recall is similar. Its additive log-odds model offers a comparatively simple probabilistic representation. Retaining all predictors allows jointly useful associations to contribute even when individual filter scores are modest; its default regularisation also limits coefficient magnitude. These properties plausibly explain its competitiveness, rather than establish the mechanism. Comparisons between PA1 and PA2 reflect both classifier and representation choices, so they do not isolate an effect of probabilistic modelling alone.

Aggressive filtering can discard complementary information. PA2's univariate ranking assesses predictors individually, while CFS favours relevance with low redundancy rather than optimising classifier F1. CFS therefore produces compact subsets but can omit information useful jointly or under a particular model representation. QDA's univariate improvement is consistent with reduced covariance-estimation demands, although it remains below Logistic Regression. Even PA2's classifier-specific wrapper does not recover the strongest complete-feature F1: greedy search bounded at 15 original predictors need not identify the most useful joint representation. PA1's counts concern transformed features, so dimensionality reductions are not directly interchangeable across applications.

Mixed Naive Bayes detects more minority-class cases than the other PA2 configurations, at the cost of lower precision and higher probability-vector error. Its conditional-independence assumption may be restrictive for related career and compensation predictors. Multiplying dependent evidence can distort estimated probabilities, although these experiments do not establish this as the sole explanation. Higher recall therefore does not imply better probability quality. AdaBoost's competitive AUC but much larger Brier similarly distinguishes ranking discrimination from accurate probability estimates.

Stacking adds an internal fitting stage and a final estimator, yet improves F1 and Brier only marginally relative to Logistic Regression, with unchanged recall and slightly lower AUC. Its related linear base learners may provide limited complementary information. Without uncertainty estimates or a formal paired comparison, these small differences do not justify claiming a meaningful advantage. Full-feature Logistic Regression is consequently competitive when balancing discrimination, probability quality and simplicity; Mixed Naive Bayes is preferable among the tested PA2 configurations when recall is the principal objective. Compact CFS subsets serve a different objective: reducing the number of recorded characteristics. No configuration dominates all objectives.

\subsection{Interpretation of Employee Characteristics}
Four predictors are selected in at least eight folds by both PA2 filters: \texttt{OverTime}, \texttt{JobLevel}, \texttt{StockOptionLevel} and \texttt{YearsAtCompany}. Overtime describes working conditions; job level describes organisational position; stock-option level represents a compensation attribute; and company tenure describes experience within the organisation. Their recurrence indicates that these complementary aspects repeatedly support predictive discrimination under different selection criteria. It does not establish that any attribute causes attrition, nor does selection alone reveal the direction or magnitude of an association. Overtime and stock-option level also appear in all three PA1 descriptive selections.

PA2's stable univariate selections further retain age, job role, involvement, monthly income and career-experience variables, supporting a multidimensional predictive interpretation. Relationships between professional level, compensation and experience also caution against interpreting each selected variable independently. PA1 agreement counts below come from separate full-data fits, whereas PA2 frequencies count outer training-fold selections. The latter indicate recurrence under overlapping resamples, not independent replication or a causal importance ranking.
The repeated original variables describe several complementary aspects of the employee context. \texttt{BusinessTravel} and \texttt{OverTime} concern working conditions; \texttt{JobInvolvement} records involvement in work; and \texttt{JobRole} distinguishes professional functions. \texttt{MonthlyIncome} and \texttt{StockOptionLevel} describe compensation, while \texttt{MaritalStatus} describes personal circumstances. Their selection by all three strategies identifies recurring associations relevant to classification in this dataset, without establishing causal effects or the direction of each association.

The existing descriptive analysis provides context for overtime: attrition is 30.53\% among employees with overtime and 10.44\% among those without it. These rates support interpreting overtime as associated with attrition-related classification patterns, rather than as a cause of leaving. \texttt{Age}, \texttt{JobSatisfaction}, \texttt{NumCompaniesWorked} and \texttt{TotalWorkingYears}, each selected by two strategies, add demographic, satisfaction and career-experience information.

Dummy categories must be interpreted together at original-variable level. For example, \texttt{OverTime\_Yes} and \texttt{OverTime\_No} encode the same employee characteristic; different retained \texttt{JobRole} categories still refer to job role. Agreement across methods makes the selected representation easier to discuss, but does not by itself provide a ranked importance measure or explain individual Neural Network predictions.

\subsection{Limitations}
The IBM dataset is synthetic/fictitious, so these findings should not be treated as evidence about a specific real workforce. Class imbalance remains present and minority-class recall is limited. Hyperparameters and decision thresholds were not optimised; conclusions concern the chosen configurations. PA1 fixes reduced subsets at 15 transformed features; PA2 ranks or searches original predictors, retaining at most 15. Neither setup establishes a globally optimal subset.

Numerical nominal codes in standalone LDA/QDA introduce representation dependence without meaningful ordering. QDA's instability inside stacking illustrates the sensitivity of covariance estimation to the available representation and sample size. PA1's wrapper and PA2's stacking fit preprocessing before their internal folds, as disclosed in Methodology; outer held-out folds remain excluded from all fitting. Cross-validation provides honest estimates for fixed configurations, but choosing the highest observed mean among many configurations may favour sampling variation. No independent external validation or statistical significance claim is provided. Logistic Regression's form is comparatively interpretable, but no fitted coefficient analysis is reported here; selection describes retained information rather than explaining individual predictions.
